\documentclass[a4paper,11pt,oneside]{book}

%%% το αρχείο αυτό καθορίζει το look που έχουν οι pythonies
%%% γίνεται \input από όλα τα κεφάλαια και τα φύλλα εργασίας

% χρησιμοποιούμενα πακέτα: 
% 
% polyglossia
% xstring
% graphicx
% caption
% xcolor
% hyperref
% minted
% geometry
% titlesec
% datetime
% changepage
% ntheorem

% οριζόμενες εντολές:
%
% smallcaps (βοηθ. removeaccents)
%    μικρά κεφαλαία χωρίς τόνους στα φωνήεντα
% scaling
%    η κλιμάκωση *όλων* των illustrations, τρέχουσα τιμή 0.9
% iconcomputer, iconkeyboard, icondiscuss, iconfillin, iconcaution, iconprompt, dottedline
%    εικονίδια για τα φύλλα εργασίας και εστιγμένη γραμμή
% marginnote
%    πλευρικό σχόλιο
% chapterwabstract (βοηθ. abstract, boxcolor, chaptercolor, concepts, tmpconcepts)
%    εισαγωγικό κείμενο κεφαλαίου με χρωματιστό τετράγωνο, συνοδευτικές έννοιες, κλπ.
% tobecontinued
%    εμφανίζει το "συνεχίζεται στην επόμενη σελίδα"

% οριζόμενα περιβάλλοντα:
% 
% note
%    μια υποσημείωση ή υπόδειξη, με μικρότερα γράμματα
% question
%    μια ερώτηση που "οδηγεί" κάθε νέα ενότητα
% answer
%    μια απάντηση σε μια ερώτηση του φύλλου εργασίας
% theory
%    μια ενότητα "θεωρίας" (στο τέλος ενός κεφαλαίου)
% exercise
%    μια αριθμημένη άσκηση
% step
%    ένα αριθμημένο βήμα (για φύλλο εργασίας)

% για μορφοποίηση κώδικα:
%
% pycode (περιβάλλον)
%     κώδικας python χωρίς αρίθμηση
% pyfile (εντολή)
%     εισαγωγή κώδικα python από αρχείο
% pyfilenl (εντολή)
%     εισαγωγή κώδικα python από αρχείο χωρίς αρίθμηση γραμμών
% pyfilesrc (εντολή)
%    εισαγωγή κώδικα από αρχείο με link στο αρχείο
% pyinline (εντολή)
%     κώδικας python μέσα στη ροή του κειμένου
% pyplain (περιβάλλον, για τα φύλλα εργασίας)
%     κώδικας python χωρίς φόντο
% pynew (περιβάλλον, για τα φύλλα εργασίας)
%     κώδικας python με φόντο
% pyterm (περιβάλλον για τα φύλλα εργασίας)
%     η είσοδος του χρήστη ή τα περιεχόμενα της οθόνης
% pyhighlight (εντολή)
%    highlight κειμένου (χρησιμοποιείται για κώδικα μέσα σε pyplain)


%%% επιλογές γλώσσας και γραμματοσειρών για το XeLaTeX

\usepackage{polyglossia}
\setdefaultlanguage{greek}
\setmainfont[Ligatures=TeX,SmallCapsFont={Linux Libertine O C},SmallCapsFeatures={Letters=SmallCaps}]{Linux Libertine O}
\setsansfont{Linux Biolinum O}
\setmonofont{Ubuntu Mono}
\enablehyphenation

% αφαίρεση τόνων από τα smallcaps
\usepackage{xstring}
\newcommand{\removeaccents}[1]{%
\def\result{#1}%
\StrSubstitute{\result}{ά}{α}[\result]%
\StrSubstitute{\result}{έ}{ε}[\result]%
\StrSubstitute{\result}{ή}{η}[\result]%
\StrSubstitute{\result}{ί}{ι}[\result]%
\StrSubstitute{\result}{ό}{ο}[\result]%
\StrSubstitute{\result}{ύ}{υ}[\result]%
\StrSubstitute{\result}{ώ}{ω}[\result]%
\StrSubstitute{\result}{Ά}{Α}[\result]%
\StrSubstitute{\result}{Έ}{Ε}[\result]%
\StrSubstitute{\result}{Ή}{Η}[\result]%
\StrSubstitute{\result}{Ί}{Ι}[\result]%
\StrSubstitute{\result}{Ό}{Ο}[\result]%
\StrSubstitute{\result}{Ύ}{Υ}[\result]%
\StrSubstitute{\result}{Ώ}{Ω}[\result]%
\result
}

\newcommand{\smallcaps}[1]{\textsc{\removeaccents{#1}}}

%%% εικόνες και λεζάντες

\usepackage{graphicx}
\newcommand{\scaling}{0.9}
\usepackage{caption}
\captionsetup{font=footnotesize}

%%% ειδικά περιβάλλοντα

\usepackage{xcolor}

% ερωτήσεις (που οδηγούν στην επόμενη ενότητα)
\definecolor{questioncolor}{rgb}{0.6,0.5,0.5}
\newenvironment{question}{\noindent\itshape\color{questioncolor}}{\noindent\ignorespaces}

% απαντήσεις (για τις ερωτήσεις των φύλλων εργασίας)
\definecolor{answercolor}{rgb}{0.5,0.5,0.5}
\newenvironment{answer}{\marginnote[16pt]{\iconfillin}\noindent\itshape\color{answercolor}}{\noindent\ignorespaces}

% περιβάλλον "θεωρίας" (πλήρες πλάτος κειμένου)
\usepackage{changepage}
\newenvironment{theory}[1]{\begin{adjustwidth}{}{-\overhang}\smallcaps{#1}\itshape}{\end{adjustwidth}}

% απομεινάρια...
% \newlength{\theoryrulelength}
% \setlength{\theoryrulelength}{36pt}
% \newenvironment{theory}{\rule{\theoryrulelength}{0.4pt}\begin{adjustwidth}{}{-\overhang}\itshape}{\end{adjustwidth}\rule{\theoryrulelength}{0.4pt}}

%%% υπερσύνδεσμοι

\definecolor{linkcolor}{rgb}{0.0,0.5,0.25}
\usepackage[colorlinks=true,urlcolor=linkcolor]{hyperref}

%%% εικονίδια και εστιγμένες γραμμές (για τα φύλλα εργασίας)

\newcommand{\iconcomputer}{\includegraphics[scale=0.35]{../../share/circle-icons/one-color/computer.eps}}
\newcommand{\iconkeyboard}{\includegraphics[scale=0.35]{../../share/circle-icons/one-color/keyboard.eps}}
\newcommand{\icondiscuss}{\includegraphics[scale=0.35]{../../share/circle-icons/one-color/chat.eps}}
\newcommand{\iconfillin}{\includegraphics[scale=0.35]{../../share/circle-icons/one-color/compose.eps}}
\newcommand{\iconcaution}{\includegraphics[scale=0.35]{../../share/circle-icons/one-color/caution.eps}}
\newcommand{\iconprompt}{\includegraphics[scale=0.35]{../../share/circle-icons/one-color/prompt.eps}}
\newcommand{\dottedline}{\vspace{\parskip}\dotfill}

%%% συνεχίζεται στην επόμενη σελίδα

\newcommand{\tobecontinued}{\mbox{}\hfill{\footnotesize ...συνεχίζεται στην επόμενη σελίδα.}}
\newenvironment{note}{\small\upshape}{}

%%% μορφοποίηση κώδικα με το pygmentize

\usepackage{minted}

% fix για ένα bug στο minted που εμφανίζεται όταν χρησιμοποιείται χρώμα στο φόντο (bgcolor)
% http://tex.stackexchange.com/questions/228058/how-to-space-before-and-after-a-minted-code-block-with-bgcolor
\makeatletter
\patchcmd{\minted@colorbg}{\noindent}{\noindent}{}{}
\apptocmd{\endminted@colorbg}{}{}{}
\makeatother

% χρώματα φόντου για τον κώδικα
\definecolor{codebg}{rgb}{0.80,0.95,0.85}
\definecolor{newcodebg}{rgb}{0.75,0.95,0.85}

% ορισμοί για τα περιβάλλοντα κώδικα
% pycode: περιβάλλον κώδικα python χωρίς αρίθμηση
\newminted[pycode]{python3}{bgcolor=codebg}
% pyfile: python από αρχείο
\newmintedfile[pyfile]{python3}{linenos=true,numberblanklines=false,escapeinside=||,bgcolor=codebg}
% pyfilenl: python από αρχείο χωρίς αρίθμηση γραμμών
\newmintedfile[pyfilenl]{python3}{linenos=false,numberblanklines=false,escapeinside=||,bgcolor=codebg}
% pyinline: python μέσα στη ροή του κειμένου
\newmintinline[pyinline]{python3}{linenos=true,numberblanklines=false}
% pyplain: (για τα φύλλα εργασίας) περιβάλλον χωρίς φόντο
\newminted[pyplain]{python3}{bgcolor=white,escapeinside=||,formatcom={\upshape}}
% pynew: (για τα φύλλα εργασίας) περιβάλλον με φόντο
\newminted[pynew]{python3}{bgcolor=newcodebg,escapeinside=||,formatcom={\upshape}}
% pyterm: (για τα φύλλα εργασίας) περιβάλλον για τα περιεχόμενα της οθόνης
\newminted[pyterm]{text}{bgcolor=white,escapeinside=||}

%\newminted[pyterm]{text}{escapeinside=||}
% [TODO] fix: το pyterm χωρίς bgcolor εμφανίζει μεγαλύτερα περιθώρια (πάνω και κάτω) και δεν φαίνεται ωραίο. Το bgcolor είναι προσωρινό workaround, έχει κι αυτό margins (για να μην είναι κολλητά ο κώδικας με το περιθώριο) κι έτσι ο κώδικας στ' αριστερά δεν είναι τέλεια στοιχισμένος.

% εντολή για κώδικα από αρχείο με link στο αρχείο
\newcommand{\pyfilesrc}[2][]{%
\pyfile[#1]{#2}\\
\mbox{}\hfill{\scriptsize\href{http://pythonies.mysch.gr/#2}{\url{#2}}}
}

% εντολή για το highlighting του κώδικα (συνήθως σε pyplain περιβάλλον με escapeinside)
\newcommand{\pyhighlight}[1]{\colorbox{newcodebg}{#1}}

%%% αριθμημένα περιβάλλοντα

\usepackage{ntheorem}

% άσκηση
\makeatletter
\theoremheaderfont{\upshape}%\upshape\bfseries\scshape}
\theorembodyfont{\itshape}%\slshape}
\newtheoremstyle{lmargin}%
  {\item[\theorem@headerfont \llap{##2}\hskip\labelsep\hskip-6pt]}%
  {\item[\theorem@headerfont \llap{##2}\hskip\labelsep ##1\ (##3)\theorem@separator]}
\makeatother
\theoremstyle{lmargin}
\newtheorem{exercise}{}[chapter]

% βήμα φύλλου εργασίας
\makeatletter
\theoremheaderfont{\bfseries}%\upshape\bfseries\scshape}
\theorembodyfont{\upshape}%\slshape}
\newtheoremstyle{lmarginup}%
  {\item[\theorem@headerfont \llap{##2}\hskip\labelsep\hskip-6pt]}%
  {\item[\theorem@headerfont \llap{##2}\hskip\labelsep ##1\ (##3)\theorem@separator]}
\newtheoremstyle{slmarginup}%
  {\item[\theorem@headerfont \llap{##1##2.}\hskip\labelsep\hskip-6pt]}%
  {\item[\theorem@headerfont \llap{##2.}\hskip\labelsep ##1\ (##3)\theorem@separator]}
\makeatother

% deprecated: \newcommand{\standalone}{} to define standalone
%\ifdefined\standalone
    \theoremstyle{slmarginup}
    \newtheorem{step}{}
%\else
%    \theoremstyle{lmarginup}
%    \newtheorem{step}{}[chapter]
%\fi

%%% γεωμετρία σελίδας και συναφείς ορισμοί από το tufte-latex
%%% https://tufte-latex.github.io/tufte-latex/

% εσοχή και διάστημα μεταξύ παραγράφων
% δεν επηρρεάζει το tufte-latex
\parindent=0pt
\parskip=6pt

% γεωμετρία σελίδας και ορισμός μηκών
\usepackage[a4paper,left=24.8mm,top=27.4mm,headsep=2\baselineskip,textwidth=107mm,marginparsep=8.2mm,marginparwidth=49.4mm,textheight=66\baselineskip,headheight=\baselineskip]{geometry}

\setlength{\marginparpush}{12pt}
\addtolength{\marginparpush}{\parskip}
\newlength{\fullwidth}
\setlength{\fullwidth}{\textwidth}
\addtolength{\fullwidth}{\marginparsep}
\addtolength{\fullwidth}{\marginparwidth}
\newlength{\overhang}
\setlength{\overhang}{\marginparsep}
\addtolength{\overhang}{\marginparwidth}

% απομεινάρια...
%\setlength\abovedisplayskip{6pt plus 2pt minus 4pt}
%\setlength\belowdisplayskip{6pt plus 2pt minus 4pt}

% italicize description run-in headings (instead of the default bold)
\renewcommand*\descriptionlabel[1]{\hspace\labelsep\normalfont\em #1}

% πλευρική σημείωση
\newcommand\marginnote[2][0pt]{%
  \marginpar{\hbox{}\vspace*{#1}\vspace*{-1\baselineskip}\noindent \footnotesize\textup{#2}}%
  {}%
}

% formatting title sections
\setcounter{secnumdepth}{-1}

\usepackage{titlesec}
\usepackage[nodate]{datetime}
\newlength{\beforesection}
\setlength{\beforesection}{3ex plus 0.5ex minus 0.2ex}
\addtolength{\beforesection}{-\parskip}
\newlength{\aftersection}
\setlength{\aftersection}{1.5ex plus 0.2ex}
\addtolength{\aftersection}{-\parskip}
\titlespacing*{\section}{0pt}{\beforesection}{\aftersection}

% απομεινάρια...
%\titlespacing*{\chapter}{0pt}{50pt}{40pt}
%\titlespacing*{\section}{0pt}{3.5ex plus 1ex minus .2ex}{2.3ex plus .2ex}

%%% για εισαγωγικό κείμενο κεφαλαίου με χρωματιστό τετράγωνο, συνοδευτικές έννοιες, κλπ.

\newcommand{\abstract}{}
\newcommand{\boxcolor}{}
\newcommand{\chaptercolor}{}
\newcommand{\concepts}{}
\newcommand{\tmpconcepts}{}
\newif\ifbonus

% reference: \titleformat{ command }[ shape ]{ format }{ label }{ sep }{ before-code }[ after-code ]
\titleformat{\chapter}[block]
{\Huge\sffamily}
{}
{0pt}
{\ifbonus\marginnote[-6pt]{\fcolorbox{\boxcolor}{\chaptercolor}{\makebox(40,40){\strut\textcolor{\boxcolor}{\Huge\thechapter}}}\\\vspace{\parskip}\\\tiny\today\\ \currenttime}\else\marginnote[-6pt]{\colorbox{\boxcolor}{\makebox(40,40){\strut\textcolor{\chaptercolor}{\Huge\thechapter}}}\\\vspace{\parskip}\\\tiny\today\\ \currenttime}\fi}
[\small\rmfamily\textmd\abstract\vspace{\parskip}\concepts\vspace{\parskip}\\\mbox{}\hrulefill]

\newcommand{\chapterwabstract}[5]{
	\renewcommand{\abstract}{#2}
    \renewcommand{\tmpconcepts}{#3}
	\ifdefempty{\tmpconcepts}{\renewcommand{\concepts}{}}{\renewcommand{\concepts}{\\\textbf{Έννοιες: }\tmpconcepts}}
	\renewcommand{\boxcolor}{#4}
	\renewcommand{\chaptercolor}{#5}
	\chapter{#1}
}

\definecolor{introColor}{rgb}{0.25,0.5,0.75}
\definecolor{answerColor}{rgb}{0.25,0.75,0.5}
\definecolor{crapsColor}{rgb}{0.5,0.75,0.25}
\definecolor{subtractionColor}{rgb}{0.5,0.25,0.75}
\definecolor{guessColor}{rgb}{0.75,0.25,0.5}
\definecolor{nimColor}{rgb}{0.75,0.5,0.25}
\definecolor{planetColor}{rgb}{0.25,0.25,0.75}
\definecolor{hangmanColor}{rgb}{0.25,0.75,0.25}
\definecolor{oxoColor}{rgb}{0.75,0.25,0.25}

\setcounter{chapter}{0}

%%% DOCUMENT START

\begin{document}

\chapterwabstract{Οδηγίες Χρήσης}{%
%Δεν είμαστε πολύ σίγουροι αν υπάρχει κανένας που διαβάζει αυτές τις εισαγωγές. Ίσως μόνο αν διστάζει να προχωρήσει στα επόμενα κεφάλαια. Γι' αυτό %θα περιοριστούμε στα απολύτως απαραίτητα.%
%Στη σύντομη αυτή εισαγωγή θα περιγράψουμε το σκεπτικό που ακολουθεί αυτό το βιβλίο, σε ποιους απευθύνεται και πως μπορεί να χρησιμοποιηθεί.
}{}{introColor}{white}

Το αντικείμενο αυτού του βιβλίου είναι ο προγραμματισμός και η ανάπτυξη ενός συγκεκριμένου \emph{τρόπου σκέψης} που τον συνοδεύει. 

Τα περισσότερα βιβλία αυτού του είδους είναι δομημένα γύρω από προγραμματιστικές έννοιες, οι οποίες εισάγονται σταδιακά. Η κατανόηση των εννοιών και η εμβάθυνση σε αυτές επιτυγχάνονται με ασκήσεις, κατά κανόνα μικρές και εστιασμένες.

Εδώ δοκιμάζουμε μια διαφορετική προσέγγιση. Κάθε κεφάλαιο αυτού του βιβλίου αντιστοιχεί σ' ένα πρόγραμμα, συνήθως ένα παιχνίδι. Δεν το παρουσιάζουμε ολοκληρωμένο για να εξηγήσουμε στη συνέχεια τον κώδικά του, αλλά ξεκινάμε από το μηδέν και το χτίζουμε βήμα προς βήμα, εξηγώντας τη φιλοσοφία με την οποία περνάμε από το κάθε βήμα στο επόμενο. Οι προγραμματιστικές έννοιες δεν είναι το επίκεντρο· εισάγονται όταν προκύψει η ανάγκη και εξηγούνται στο βαθμό που εξυπηρετεί το συγκεκριμένο παράδειγμα.

Πιστεύουμε ότι η κατασκευή μεγαλύτερων προγραμμάτων παρέχει στους αρχάριους βαθύτερο κίνητρο και αμεσότερη ανταμοιβή. Ελπίζουμε ότι ο τρόπος ανάπτυξης του περιεχομένου επιτρέπει σ' εκείνους που το χρησιμοποιούν να γενικεύουν και να εφαρμόζουν τις γνώσεις που αποκτούν και σε νέα προβλήματα.

Τα προγράμματα που κατασκευάζουμε δεν έχουν γραφικά ή ήχο, ώστε να μπορέσουμε να εστιάσουμε στα θεμελιώδη. Κάποια από τα κεφάλαια, όπως το \emph{Μάντεψε τον Αριθμό} και η \emph{Κρεμάλα}, αφορούν κλασικά παραδείγματα που χρησιμοποιούνται παραδοσιακά σε εισαγωγικά μαθήματα επειδή είναι αρκετά απλά αλλά και ιδιαίτερα περιεκτικά. Aυτά έχουμε προσπαθήσει να τα εμπλουτίσουμε με εναλλακτικές πινελιές. Ωστόσο, τα περισσότερα κεφάλαια ασχολούνται με λιγότερο διαδεδομένα ή και εντελώς καινούργια παραδείγματα. 
%Πηγή έμπνευσης για πολλές από τις ασκήσεις του βιβλίου αποτέλεσε η εξαιρετική συλλογή %
%\marginnote{\href{http://www.vintage-basic.net/games.html}{\url{vintage-basic.net}}}
%\emph{BASIC Computer Games} που επιμελήθηκε ο David Ahl. Δημοσιεύτηκε το 1978, μια εποχή που ήταν πολύ συνηθισμένο να πληκτρολογείς ολόκληρα προγράμματα και να προσπαθείς να καταλάβεις πως δουλεύουν.

Θα γράφουμε τα προγράμματά μας στη γλώσσα Python. 
\marginnote[2pt]{Οι εξηγήσεις που σχετίζονται ειδικά με την Python έχουν διαχωριστεί από το βασικό κείμενο και βρίσκονται εδώ, στο περιθώριο.}
Θα μπορούσαμε να διαλέξουμε μια άλλη γλώσσα, όμως η Python είναι εξαιρετική για τον σκοπό μας.
Προσπαθήσαμε το κείμενό μας να είναι ανεξάρτητο από τη γλώσσα προγραμματισμού. Οτιδήποτε έχει να κάνει ειδικά με την Python θα βρίσκεται στο \emph{περιθώριο}.%

Αν θέλετε να χρησιμοποιήσετε αυτό το βιβλίο για να μάθετε μόνοι σας προγραμματισμό, πρέπει να γνωρίζετε ότι θα χρειαστεί να επενδύσετε αρκετό χρόνο και προσπάθεια. Όπως συμβαίνει και με πολλά άλλα πράγματα, δεν γίνεται να μάθετε προγραμματισμό διαβάζοντας απλά ένα βιβλίο. Εξετάστε τον κώδικα που δίνεται, προσπαθήστε να τον κατανοήσετε διαβάζοντας το συνοδευτικό κείμενο και πληκτρολογήστε τον για να τον δείτε να εκτελείται. Αν το πρόγραμμα δεν εκτελείται ή αν δεν συμπεριφέρεται όπως περιμένετε τότε είστε τυχεροί: έχετε κάνει κάποιο λάθος και η διαδικασία που θα ακολουθήσετε για να το εντοπίσετε και να το διορθώσετε είναι αυτή που θα σας προσφέρει μια βαθύτερη κατανόηση.

Καθώς πληκτρολογείτε τον κώδικα, δώστε προσοχή στους αριθμούς των γραμμών που βρίσκονται αριστερά. Αυτοί θα σας βοηθήσουν να καταλάβετε σε ποιο σημείο θα εισάγετε νέο κώδικα ή θα αντικαταστήσετε τμήματα του ήδη υπάρχοντος. Αν ο κώδικας δεν είναι αριθμημένος, τότε πρόκειται απλά για μια εναλλακτική πρόταση.

Πιθανώς να θέλετε να χρησιμοποιήσετε αυτό το βιβλίο για να διδάξετε προγραμματισμό. Κι εμείς έτσι το χρησιμοποιούμε. Ίσως η προσέγγισή του ταιριάζει περισσότερο σε σας ή τους μαθητές σας. Σε αυτή την περίπτωση, θεωρήστε ότι το κείμενο που συνοδεύει τον κώδικα αποτελεί μια πρόταση ως προς τα σημεία στα οποία θα πρέπει να εστιάσετε τη διδασκαλία σας. 
% Προσπαθήσαμε σε κάθε περίπτωση να χρησιμοποιήσουμε κώδικα που είναι κατανοητός από τους μαθητές
% Δεν αποκλείεται σε ορισμένες περιπτώσεις να σκεφτείτε ότι εσείς θα γράφατε διαφορετικά ένα κομμάτι κώδικα. 

Ξεκινώντας, είναι απαραίτητο να %
\marginnote{Μπορείτε να κατεβάσετε την Python από εδώ: \href{https://www.python.org/downloads/}{\url{python.org/downloads}}}%
κατεβάσετε και να εγκαταστήσετε την \href{https://www.python.org/downloads/}{Python}. %
\marginnote{Η Python κυκλοφορεί σε δύο εκδόσεις που δεν είναι συμβατές μεταξύ τους. Το βιβλίο ακολουθεί την \emph{Python 3}.} Σε περίπτωση που χρειάζεστε βοήθεια, ανατρέξτε στην %
\marginnote{Στη διεύθυνση \href{http://pythonies.mysch.gr/}{\url{pythonies.mysch.gr}} θα βρείτε όλο το υλικό του βιβλίου.}%
\href{http://pythonies.mysch.gr/}{ιστοσελίδα} μας για κάποιες βασικές οδηγίες, συμβουλευτείτε οποιαδήποτε από τις \href{http://users.sch.gr/boukeas/sources/}{πηγές} που προτείνουμε στη βιβλιογραφία, χρησιμοποιήστε τον %
\marginnote{Ο «επίσημος» εισαγωγικός οδηγός στη διεύθυνση: \href{https://wiki.python.org/moin/BeginnersGuide}{\url{wiki.python.org/moin/BeginnersGuide}}}%
\href{https://wiki.python.org/moin/BeginnersGuide}{οδηγό για αρχαρίους} (στ' Αγγλικά) ή κάντε μια αναζήτηση στον Παγκόσμιο Ιστό.

%Ένα άλλο βιβλίο που αποτέλεσε πηγή έμπνευσης είναι \emph{BASIC Computer Games}, του David Ahl. Περιέχει κώδικα για δεκάδες παιχνίδια κειμένου, γραμμένα στη γλώσσα BASIC. Ο σχολιασμός που συνοδεύει τα προγράμματα είναι σχεδόν ανύπαρκτος. Γράφτηκε το 1978, μια εποχή που ήταν πολύ πιο συνηθισμένο να πληκτρολογείς ολόκληρα προγράμματα και να προσπαθείς να καταλάβεις πως δουλεύουν.
%http://www.vintage-basic.net/games.html

\end{document}

\section{Βιβλιογραφία}

% πρόσθεσε το problem space

Αυτό το βιβλίο δεν είναι το μοναδικό διαθέσιμο αν θέλετε να μάθετε προγραμματισμό σε Python. 

Αν είστε άνετοι με τα Αγγλικά, η σειρά βιβλίων του Al Sweigart και ειδικότερα το %
\marginnote{\href{https://inventwithpython.com}{\url{inventwithpython.com}}}%
\emph{Invent Your Own Computer Games in Python} είναι μια εξαιρετική επιλογή. 
Η προσέγγιση και εδώ βασίζεται σε ολοκληρωμένα παραδείγματα παιχνιδιών.

Δύο ακόμα καλές επιλογές στα Αγγλικά είναι το %
\marginnote{\href{http://greenteapress.com/thinkpython}{\url{greenteapress.com/thinkpython}}}%
\emph{How to Think Like a Computer Scientist} του Allen Downey και το %
\marginnote{\href{http://learnpythonthehardway.org/book}{\url{learnpythonthehardway.org/book}}}%
\emph{Learn Python the Hard Way} του Zed Shaw.

Μπορείτε βέβαια πάντα να ακολουθήσετε τα %
\marginnote{\href{https://docs.python.org/3.4/tutorial/}{\url{docs.python.org/3.4/tutorial}}}
σύντομα και περιεκτικά παραδείγματα στην ιστοσελίδα τεκμηρίωσης της Python.

Τελευταία έχουν αρχίσει να εμφανίζονται βιβλία για την Python και στα ελληνικά. Το %
\marginnote{\href{https://pygamegr.wordpress.com}{\url{pygamegr.wordpress.com}}}%
\emph{Παιχνίδια σε Python \& Pygame} του Μανώλη Κιαγιά αφορά παιχνίδια με γραφικά και είναι πολύ ευχάριστο, ενώ το %
\marginnote{\href{http://aggelid.mysch.gr/pythonbook}{\url{aggelid.mysch.gr/pythonbook}}}%
\emph{Εισαγωγή στον Προγραμματισμό με την Python} του Νίκου Αγγελιδάκη είναι πληρέστατο και ακολουθεί μια πιο παραδοσιακή προσέγγιση.




